\begin{abstract}
Disorders of the visual pathway alter both anatomy, visible on structural MRI, and signal timing, recorded as the visual evoked potential (VEP), yet most computational tools use only one of the two, and no public dataset pairs them with diagnosis labels. We present the Cross-Attention Fusion Network (CAFN), which couples a squeeze-and-excitation 3D residual encoder and a multi-scale 1D encoder through gated bidirectional cross-attention, and evaluate it in two ways. In a controlled synthetic benchmark in which the class information carried by each modality is tunable, fusion improved test accuracy by \gainUni{} percentage points over the better single modality ($\accFull\%$ versus $\accMRI\%$ and $\accVEP\%$), but cross-attention was statistically indistinguishable from concatenation ($\accCat\%$) and element-wise sum ($\accSum\%$), and none of seven component ablations produced a reliable change. On two public OpenNeuro cohorts (95 subjects; depression with MEG and zoster neuralgia with sleep EEG, each with T1-weighted MRI), evaluated with pooled and leave-one-cohort-out protocols, repeated cross-validation and permutation tests, neural fusion was at chance (AUC 0.51), a weak spectral signal was detectable (AUC 0.56--0.62), and nothing transferred between cohorts. The one-command pipeline runs from download to test statistics under a fixed seed and is supported by a 340-test suite. We report where fusion helps, where it does not, and the pitfalls we found; no diagnostic accuracy is claimed.
\end{abstract}
